\documentclass[11pt,a4paper]{article}
\usepackage[utf8]{inputenc}
\usepackage[T1]{fontenc}
\usepackage{amsmath,amssymb,amsfonts,amsthm}
\usepackage{geometry}
\usepackage{hyperref}
\usepackage{enumitem}

\geometry{margin=1in}
\allowdisplaybreaks

\title{\textbf{Complete Explicit-Index Derivations for Lecture 7--9:\\ Symmetries and Fields (Poincar\'e and Lorentz Group Representations)}}
\author{}
\date{}

\begin{document}

\maketitle

\section*{Global Conventions Established at the Outset}
\begin{enumerate}[label=\textbf{\arabic*.}]
    \item \textbf{Spacetime Metric}: Minkowski spacetime $\mathbb{R}^{1,3}$ uses the mostly-plus metric signature:
    \begin{equation}
        \eta_{\mu\nu} = \mathrm{diag}(-1, +1, +1, +1), \qquad \eta^{\mu\nu} = \mathrm{diag}(-1, +1, +1, +1),
    \end{equation}
    with Greek indices $\mu,\nu,\rho,\sigma \in \{0,1,2,3\}$ and Latin spatial indices $i,j,k \in \{1,2,3\}$. For spatial components, raising or lowering an index does not change the sign ($\eta_{ij} = \delta_{ij}$), whereas raising or lowering a $0$-index introduces a minus sign ($\eta_{00} = -1$).
    \item \textbf{Kronecker Delta}: $\delta^\mu{}_\nu = \eta^{\mu\rho}\eta_{\rho\nu}$ has value $1$ if $\mu=\nu$ and $0$ otherwise.
    \item \textbf{3D Levi-Civita Symbol}: $\epsilon_{ijk} = \epsilon^{ijk}$ is completely antisymmetric with $\epsilon_{123} = +1$. It satisfies the contraction identities:
    \begin{equation}
        \epsilon_{ijk}\epsilon_{imn} = \delta_{jm}\delta_{kn} - \delta_{jn}\delta_{km}, \qquad \epsilon_{ijk}\epsilon_{ijn} = 2\delta_{kn}.
    \end{equation}
    \item \textbf{4D Levi-Civita Tensor}: We choose the convention $\epsilon^{0123} = -1$ (so $\epsilon_{0123} = +1$) so that $\epsilon^{0ijk} = -\epsilon_{ijk}$.
\end{enumerate}

\vspace{0.5em}\hrule\vspace{1em}

\section{Lorentz Group and the $\mathfrak{so}(1,3)$ Lie Algebra (Section 4.1.1)}

\subsection{Vector Representation and Infinitesimal Generators (Eqs. 4.1 -- 4.9)}
A Lorentz transformation acts linearly on spacetime coordinates $x^\mu$ as:
\begin{equation}
    x'^\mu = \Lambda^\mu{}_\nu x^\nu.
\end{equation}
Preservation of the Minkowski inner product $x'^\mu \eta_{\mu\nu} x'^\nu = x^\rho \eta_{\rho\sigma} x^\sigma$ for arbitrary $x$ gives the defining condition of $O(1,3)$ in explicit index notation (avoiding the matrix shorthand $\Lambda^T \eta \Lambda = \eta$ in Eq.~4.6):
\begin{equation}
    \Lambda^\mu{}_\rho \eta_{\mu\nu} \Lambda^\nu{}_\sigma = \eta_{\rho\sigma}.
\end{equation}
Multiplying both sides by $\eta^{\sigma\tau} (\Lambda^{-1})^\lambda{}_\rho$ yields the inverse transformation rule:
\begin{equation}
    (\Lambda^{-1})^\tau{}_\nu = \eta^{\tau\sigma}\eta_{\nu\mu}\Lambda^\mu{}_\sigma = \Lambda_\nu{}^\tau.
\end{equation}

\paragraph{Convention Choice (Matrix Entries of $l^{\mu\nu}$):}
For each antisymmetric pair of indices $[\mu\nu]$, the generator $l^{\mu\nu} = -l^{\nu\mu}$ is a linear operator on $\mathbb{R}^{1,3}$ whose components are $(l^{\mu\nu})^\alpha{}_\beta$, where $\alpha$ is the row (upper) index and $\beta$ is the column (lower) index. In explicit index notation, the exponential map (Eq.~4.7) is:
\begin{equation}
    \Lambda^\alpha{}_\beta = \left(\exp\left[-\frac{i}{2}\omega_{\mu\nu}l^{\mu\nu}\right]\right)^\alpha{}_\beta = \delta^\alpha{}_\beta - \frac{i}{2}\omega_{\mu\nu}(l^{\mu\nu})^\alpha{}_\beta + \mathcal{O}(\omega^2),
\end{equation}
where $\omega_{\mu\nu} = -\omega_{\nu\mu}$. Substituting this into the invariance condition $\Lambda^\gamma{}_\alpha \eta_{\gamma\delta} \Lambda^\delta{}_\beta = \eta_{\alpha\beta}$ to first order in $\omega$:
\begin{align}
    \eta_{\alpha\beta} &= \left(\delta^\gamma{}_\alpha - \frac{i}{2}\omega_{\mu\nu}(l^{\mu\nu})^\gamma{}_\alpha\right)\eta_{\gamma\delta}\left(\delta^\delta{}_\beta - \frac{i}{2}\omega_{\mu\nu}(l^{\mu\nu})^\delta{}_\beta\right) + \mathcal{O}(\omega^2) \nonumber \\
    &= \eta_{\alpha\beta} - \frac{i}{2}\omega_{\mu\nu}\left[\eta_{\gamma\beta}(l^{\mu\nu})^\gamma{}_\alpha + \eta_{\alpha\delta}(l^{\mu\nu})^\delta{}_\beta\right] + \mathcal{O}(\omega^2).
\end{align}
Defining the fully lowered generator matrix $(l^{\mu\nu})_{\alpha\beta} \equiv \eta_{\alpha\gamma}(l^{\mu\nu})^\gamma{}_\beta$, this requires $(l^{\mu\nu})_{\alpha\beta} = -(l^{\mu\nu})_{\beta\alpha}$. To match the Poincar\'e algebra (4.4) and the explicit matrices (4.9), the unique choice of basis generators is:
\begin{equation}
    (l^{\mu\nu})^\alpha{}_\beta = i\left(\eta^{\mu\alpha}\delta^\nu{}_\beta - \eta^{\nu\alpha}\delta^\mu{}_\beta\right).
\end{equation}

\paragraph{Verification of Eq. (4.9):}
\begin{itemize}
    \item For $\mu=0, \nu=1$:
    \begin{equation}
        (l^{01})^\alpha{}_\beta = i\left(\eta^{0\alpha}\delta^1{}_\beta - \eta^{1\alpha}\delta^0{}_\beta\right).
    \end{equation}
    Since $\eta^{00} = -1$ and $\eta^{11} = +1$:
    \begin{equation}
        (l^{01})^0{}_1 = i(-1)(1) = -i, \qquad (l^{01})^1{}_0 = i(0 - (1)(1)) = -i.
    \end{equation}
    \emph{Typo Alert in Eq. (4.9):} The lecture notes write $l^{01} = +i\begin{pmatrix} 0 & 1 & 0 & 0 \\ 1 & 0 & 0 & 0 \\ 0 & 0 & 0 & 0 \\ 0 & 0 & 0 & 0 \end{pmatrix}$. However, with $\eta^{00}=-1$, $(l^{01})^\alpha{}_\beta$ satisfying the Lorentz algebra (4.4) with $\eta_{\mu\nu}=\mathrm{diag}(-1,1,1,1)$ must have $-i$ in the $(0,1)$ and $(1,0)$ entries (or equivalently, $+i$ corresponds to $l_{01} = \eta_{00}\eta_{11}l^{01} = -l^{01}$).
    \item For $\mu=1, \nu=2$:
    \begin{equation}
        (l^{12})^\alpha{}_\beta = i\left(\eta^{1\alpha}\delta^2{}_\beta - \eta^{2\alpha}\delta^1{}_\beta\right) \implies (l^{12})^1{}_2 = +i = -i(-1), \quad (l^{12})^2{}_1 = -i = +i(-1),
    \end{equation}
    which gives:
    \begin{equation}
        (l^{12})^\alpha{}_\beta = -i\begin{pmatrix} 0 & 0 & 0 & 0 \\ 0 & 0 & -1 & 0 \\ 0 & 1 & 0 & 0 \\ 0 & 0 & 0 & 0 \end{pmatrix}^\alpha{}_\beta.
    \end{equation}
\end{itemize}

\paragraph{Explicit Index Derivation of the Lorentz Commutation Relations (Eq. 4.4):}
Compute the commutator $[l^{\mu\nu}, l^{\rho\sigma}]^\alpha{}_\beta = (l^{\mu\nu})^\alpha{}_\gamma (l^{\rho\sigma})^\gamma{}_\beta - (l^{\rho\sigma})^\alpha{}_\gamma (l^{\mu\nu})^\gamma{}_\beta$:
\begin{align}
    (l^{\mu\nu})^\alpha{}_\gamma (l^{\rho\sigma})^\gamma{}_\beta &= i^2 \left(\eta^{\mu\alpha}\delta^\nu{}_\gamma - \eta^{\nu\alpha}\delta^\mu{}_\gamma\right)\left(\eta^{\rho\gamma}\delta^\sigma{}_\beta - \eta^{\sigma\gamma}\delta^\rho{}_\beta\right) \nonumber \\
    &= -\left(\eta^{\mu\alpha}\eta^{\rho\nu}\delta^\sigma{}_\beta - \eta^{\mu\alpha}\eta^{\sigma\nu}\delta^\rho{}_\beta - \eta^{\nu\alpha}\eta^{\rho\mu}\delta^\sigma{}_\beta + \eta^{\nu\alpha}\eta^{\sigma\mu}\delta^\rho{}_\beta\right).
\end{align}
Swapping $(\mu,\nu) \leftrightarrow (\rho,\sigma)$ gives the second term:
\begin{equation}
    (l^{\rho\sigma})^\alpha{}_\gamma (l^{\mu\nu})^\gamma{}_\beta = -\left(\eta^{\rho\alpha}\eta^{\mu\sigma}\delta^\nu{}_\beta - \eta^{\rho\alpha}\eta^{\nu\sigma}\delta^\mu{}_\beta - \eta^{\sigma\alpha}\eta^{\mu\rho}\delta^\nu{}_\beta + \eta^{\sigma\alpha}\eta^{\nu\rho}\delta^\mu{}_\beta\right).
\end{equation}
Subtracting the two and grouping by $\eta^{\nu\rho}, \eta^{\nu\sigma}, \eta^{\mu\sigma}, \eta^{\mu\rho}$:
\begin{align}
    [l^{\mu\nu}, l^{\rho\sigma}]^\alpha{}_\beta &= \eta^{\nu\rho}\left(\eta^{\sigma\alpha}\delta^\mu{}_\beta - \eta^{\mu\alpha}\delta^\sigma{}_\beta\right) - \eta^{\nu\sigma}\left(\eta^{\rho\alpha}\delta^\mu{}_\beta - \eta^{\mu\alpha}\delta^\rho{}_\beta\right) \nonumber \\
    &\quad + \eta^{\mu\sigma}\left(\eta^{\rho\alpha}\delta^\nu{}_\beta - \eta^{\nu\alpha}\delta^\rho{}_\beta\right) - \eta^{\mu\rho}\left(\eta^{\sigma\alpha}\delta^\nu{}_\beta - \eta^{\nu\alpha}\delta^\sigma{}_\beta\right) \nonumber \\
    &= i\left[\eta^{\nu\rho}(l^{\mu\sigma})^\alpha{}_\beta - \eta^{\nu\sigma}(l^{\mu\rho})^\alpha{}_\beta + \eta^{\mu\sigma}(l^{\nu\rho})^\alpha{}_\beta - \eta^{\mu\rho}(l^{\nu\sigma})^\alpha{}_\beta\right].
\end{align}
This reproduces Eq.~(4.4).

\paragraph{Derivation of $[L^{\mu\nu}, P^\sigma]$ (Eq. 4.3):}
In the coordinate representation on scalar functions, $P^\sigma = -i\partial^\sigma$ and $L^{\mu\nu} = x^\mu P^\nu - x^\nu P^\mu = -i(x^\mu \partial^\nu - x^\nu \partial^\mu)$. Acting on a test function $f(x)$:
\begin{align}
    [L^{\mu\nu}, P^\sigma]f &= \left[-i(x^\mu \partial^\nu - x^\nu \partial^\mu), -i\partial^\sigma\right]f \nonumber \\
    &= -\left(x^\mu \partial^\nu \partial^\sigma f - x^\nu \partial^\mu \partial^\sigma f - \partial^\sigma(x^\mu \partial^\nu f - x^\nu \partial^\mu f)\right) \nonumber \\
    &= (\partial^\sigma x^\mu)\partial^\nu f - (\partial^\sigma x^\nu)\partial^\mu f \nonumber \\
    &= \eta^{\sigma\mu}\partial^\nu f - \eta^{\sigma\nu}\partial^\mu f = i\left(\eta^{\sigma\mu}P^\nu - \eta^{\sigma\nu}P^\mu\right)f = i\left(P^\mu\eta^{\nu\sigma} - P^\nu\eta^{\mu\sigma}\right)f,
\end{align}
which matches Eq.~(4.3).

\subsection{Rotations, Boosts, and the Typo in $j_i$ (Eqs. 4.10 -- 4.12)}
\paragraph{Typo Resolution in Eq. (4.10):}
The lecture notes write $j_i = \frac{1}{2}\epsilon_{ijk}l^{jk}$ and $k_i = l^{0i}$. With $\eta_{ij} = +\delta_{ij}$ and the Lorentz algebra (4.4), $j_i$ \textbf{must have a minus sign in front}:
\begin{equation}
    j_i = -\frac{1}{2}\epsilon_{ijk}l^{jk} \iff l^{mn} = -\epsilon_{mnk}j_k, \qquad k_i = l^{0i}.
\end{equation}
Let us check the commutator $[j_i, j_j]$ using Eq.~(4.4) with $\eta^{mn} = +\delta^{mn}$:
\begin{align}
    [j_i, j_j] &= \frac{1}{4}\epsilon_{imn}\epsilon_{jpq}[l^{mn}, l^{pq}] \nonumber \\
    &= \frac{i}{4}\epsilon_{imn}\epsilon_{jpq}\left(\delta^{np}l^{mq} - \delta^{nq}l^{mp} + \delta^{mq}l^{np} - \delta^{mp}l^{nq}\right) \nonumber \\
    &= \frac{i}{4}\epsilon_{imn}\epsilon_{jpq}\left(4\delta^{np}l^{mq}\right) = i\epsilon_{imn}\epsilon_{jnq}l^{mq} \nonumber \\
    &= i\epsilon_{nim}\epsilon_{njq}\left(-\epsilon_{mqk}j_k\right) = -i\left(\delta_{ij}\delta_{mq} - \delta_{iq}\delta_{mj}\right)\epsilon_{mqk}j_k = +i\epsilon_{ijk}j_k.
\end{align}
If $j_i = +\frac{1}{2}\epsilon_{ijk}l^{jk}$ were used without the minus sign, one would obtain $[j_i, j_j] = -i\epsilon_{ijk}j_k$, violating Eq.~(4.12).

\paragraph{Hermitian Conjugation Properties (Eq. 4.11):}
In explicit index notation, the matrix adjoint is $((l^{\mu\nu})^\dagger)^\alpha{}_\beta \equiv ((l^{\mu\nu})^\beta{}_\alpha)^*$:
\begin{itemize}
    \item For spatial indices $j,k \in \{1,2,3\}$, $\eta^{j\alpha} = \delta^{j\alpha}$, so:
    \begin{align}
        \left((j_i)^\dagger\right)^\alpha{}_\beta &= \left((j_i)^\beta{}_\alpha\right)^* = -\frac{1}{2}\epsilon_{ijk}\left[i\left(\delta^{j\beta}\delta^k{}_\alpha - \delta^{k\beta}\delta^j{}_\alpha\right)\right]^* \nonumber \\
        &= +\frac{i}{2}\epsilon_{ijk}\left(\delta^{j\beta}\delta^k{}_\alpha - \delta^{k\beta}\delta^j{}_\alpha\right) = -\frac{i}{2}\epsilon_{ijk}\left(\delta^{j\alpha}\delta^k{}_\beta - \delta^{k\alpha}\delta^j{}_\beta\right) = (j_i)^\alpha{}_\beta.
    \end{align}
    Hence $j_i^\dagger = j_i$.
    \item For boost generators $k_i = l^{0i}$, the nonzero entries have $\{\alpha,\beta\} = \{0,i\}$:
    \begin{equation}
        (k_i)^0{}_i = i(\eta^{00}\delta^i{}_i - \eta^{i0}\delta^0{}_i) = -i, \qquad (k_i)^i{}_0 = i(\eta^{0i}\delta^i{}_0 - \eta^{ii}\delta^0{}_0) = -i.
    \end{equation}
    Taking the conjugate transpose:
    \begin{equation}
        \left((k_i)^\dagger\right)^0{}_i = \left((k_i)^i{}_0\right)^* = (-i)^* = +i = -(k_i)^0{}_i, \qquad \left((k_i)^\dagger\right)^i{}_0 = \left((k_i)^0{}_i\right)^* = +i = -(k_i)^i{}_0.
    \end{equation}
    Hence $k_i^\dagger = -k_i$.
\end{itemize}

\paragraph{Derivation of $[j_i, k_j]$ and $[k_i, k_j]$ (Eq. 4.12):}
Using Eq.~(4.4) with $\eta^{00} = -1$, $\eta^{0m} = 0$, and $\eta^{mn} = \delta_{mn}$:
\begin{enumerate}
    \item \textbf{Commutator $[j_i, k_j]$}:
    \begin{align}
        [j_i, k_j] &= -\frac{1}{2}\epsilon_{imn}[l^{mn}, l^{0j}] \nonumber \\
        &= -\frac{i}{2}\epsilon_{imn}\left(\eta^{n0}l^{mj} - \eta^{nj}l^{m0} + \eta^{mj}l^{n0} - \eta^{m0}l^{nj}\right) \nonumber \\
        &= -\frac{i}{2}\epsilon_{imn}\left(-\delta_{nj}(-l^{0m}) + \delta_{mj}(-l^{0n})\right) \nonumber \\
        &= -\frac{i}{2}\left(\epsilon_{imj}k_m - \epsilon_{ijn}k_n\right) = i\epsilon_{ijk}k_k.
    \end{align}
    \item \textbf{Commutator $[k_i, k_j]$}:
    \begin{align}
        [k_i, k_j] &= [l^{0i}, l^{0j}] = i\left(\eta^{i0}l^{0j} - \eta^{ij}l^{00} + \eta^{0j}l^{i0} - \eta^{00}l^{ij}\right) \nonumber \\
        &= i\left(-(-1)l^{ij}\right) = i l^{ij} = -i\epsilon_{ijk}j_k.
    \end{align}
\end{enumerate}
All three relations in Eq.~(4.12) are thus established:
\begin{equation}
    [j_i, j_j] = i\epsilon_{ijk}j_k, \qquad [j_i, k_j] = i\epsilon_{ijk}k_k, \qquad [k_i, k_j] = -i\epsilon_{ijk}j_k.
\end{equation}

\subsection{Chiral Decomposition $\mathfrak{so}(1,3)_{\mathbb{C}} \simeq \mathfrak{sl}(2,\mathbb{C})_L \oplus \mathfrak{sl}(2,\mathbb{C})_R$ (Eqs. 4.13 -- 4.24)}
Define the complexified combinations (Eq.~4.14):
\begin{equation}
    A_i = \frac{1}{2}(j_i + ik_i), \qquad B_i = \frac{1}{2}(j_i - ik_i).
\end{equation}
Their commutators (Eq.~4.15) follow directly from Eq.~(4.12):
\begin{align}
    [A_i, A_j] &= \frac{1}{4}[j_i + ik_i, j_j + ik_j] = \frac{1}{4}\left([j_i, j_j] + i[j_i, k_j] + i[k_i, j_j] - [k_i, k_j]\right) \nonumber \\
    &= \frac{1}{4}\left(i\epsilon_{ijk}j_k + i(i\epsilon_{ijk}k_k) + i(i\epsilon_{ijk}k_k) - (-i\epsilon_{ijk}j_k)\right) \nonumber \\
    &= \frac{1}{2}i\epsilon_{ijk}(j_k + ik_k) = i\epsilon_{ijk}A_k, \\[6pt]
    [B_i, B_j] &= \frac{1}{4}[j_i - ik_i, j_j - ik_j] = \frac{1}{4}\left([j_i, j_j] - i[j_i, k_j] - i[k_i, j_j] - [k_i, k_j]\right) \nonumber \\
    &= \frac{1}{4}\left(i\epsilon_{ijk}j_k - i(i\epsilon_{ijk}k_k) - i(i\epsilon_{ijk}k_k) - (-i\epsilon_{ijk}j_k)\right) \nonumber \\
    &= \frac{1}{2}i\epsilon_{ijk}(j_k - ik_k) = i\epsilon_{ijk}B_k, \\[6pt]
    [A_i, B_j] &= \frac{1}{4}[j_i + ik_i, j_j - ik_j] = \frac{1}{4}\left([j_i, j_j] - i[j_i, k_j] + i[k_i, j_j] + [k_i, k_j]\right) \nonumber \\
    &= \frac{1}{4}\left(i\epsilon_{ijk}j_k - i(i\epsilon_{ijk}k_k) + i(i\epsilon_{ijk}k_k) + (-i\epsilon_{ijk}j_k)\right) = 0.
\end{align}

\paragraph{Reality Condition (Eq. 4.21):}
Following the lecture notes' convention where complex conjugation $*$ on the complexified algebra reverses the $i$ in front of the real Lie algebra generators $(-ij_i)^* = -ij_i \implies j_i^* = -j_i$ and $(-ik_i)^* = -ik_i \implies k_i^* = -k_i$:
\begin{equation}
    A_i^* = \left(\frac{1}{2}(j_i + ik_i)\right)^* = \frac{1}{2}\left(j_i^* - ik_i^*\right) = \frac{1}{2}\left(-j_i + ik_i\right) = -B_i,
\end{equation}
which gives Eq.~(4.21), and implies that complex conjugation of a representation swaps the $(j_L, j_R)$ labels: $(j_L, j_R)^* = (j_R, j_L)$ (Eq.~4.23).

\vspace{0.5em}\hrule\vspace{1em}

\section{Spinors, $SL(2,\mathbb{C})$, and Index Conventions (Section 4.1.2)}

\subsection{Index Structure of $\epsilon_{\alpha\beta}$, $\epsilon^{\alpha\beta}$, and Spinor Raising/Lowering (Eqs. 4.47 -- 4.71)}
\paragraph{Convention Choices (Professor's Convention $\epsilon_{\dot{\alpha}\dot{\beta}} = -\epsilon_{\alpha\beta}$):}
\begin{itemize}
    \item Undotted indices $\alpha,\beta,\gamma,\delta \in \{1,2\}$ label the fundamental $(\frac{1}{2},0)$ representation of $SL(2,\mathbb{C})$.
    \item Dotted indices $\dot{\alpha},\dot{\beta},\dot{\gamma},\dot{\delta} \in \{1,2\}$ label the complex-conjugate $(0,\frac{1}{2})$ representation of $SL(2,\mathbb{C})$.
    \item \textbf{Antisymmetric Spinor Metric (Eq. 4.55 \& Professor's Dotted Convention):}
    \begin{align}
        \epsilon_{12} &= +1, \quad \epsilon_{21} = -1, \quad \epsilon_{11}=\epsilon_{22}=0; \qquad \epsilon^{12} = -1, \quad \epsilon^{21} = +1, \quad \epsilon^{11}=\epsilon^{22}=0, \\
        \epsilon_{\dot{1}\dot{2}} &= -1, \quad \epsilon_{\dot{2}\dot{1}} = +1, \quad \epsilon_{\dot{1}\dot{1}}=\epsilon_{\dot{2}\dot{2}}=0; \qquad \epsilon^{\dot{1}\dot{2}} = +1, \quad \epsilon^{\dot{2}\dot{1}} = -1, \quad \epsilon^{\dot{1}\dot{1}}=\epsilon^{\dot{2}\dot{2}}=0.
    \end{align}
    Note that because $\epsilon_{\dot{\alpha}\dot{\beta}} = -\epsilon_{\alpha\beta}$ and $\epsilon^{\dot{\alpha}\dot{\beta}} = -\epsilon^{\alpha\beta}$, the diagram of index-raising and Hermitian conjugation ($\dagger$) anticommutes: if $(\psi_\alpha)^\dagger = \bar{\psi}_{\dot{\alpha}}$, then:
    \begin{equation}
        (\psi^\alpha)^\dagger = (\epsilon^{\alpha\beta}\psi_\beta)^\dagger = \epsilon^{\alpha\beta}\bar{\psi}_{\dot{\beta}} = -\epsilon^{\dot{\alpha}\dot{\beta}}\bar{\psi}_{\dot{\beta}} = -\bar{\psi}^{\dot{\alpha}}.
    \end{equation}
\end{itemize}

\paragraph{Verification of Eq. (4.56) and (4.70):}
Contracting the second index of $\epsilon^{\alpha\beta}$ with the first index of $\epsilon_{\beta\gamma}$ (and identically for dotted indices, since both $\epsilon^{\dot{\alpha}\dot{\beta}}$ and $\epsilon_{\dot{\beta}\dot{\gamma}}$ flip sign):
\begin{equation}
    \epsilon^{\alpha\beta}\epsilon_{\beta\gamma} = \delta^\alpha{}_\gamma, \qquad \epsilon^{\alpha\beta}\epsilon_{\gamma\beta} = -\delta^\alpha{}_\gamma, \qquad \epsilon^{\dot{\alpha}\dot{\beta}}\epsilon_{\dot{\beta}\dot{\gamma}} = \delta^{\dot{\alpha}}{}_{\dot{\gamma}}, \qquad \epsilon^{\dot{\alpha}\dot{\beta}}\epsilon_{\dot{\gamma}\dot{\beta}} = -\delta^{\dot{\alpha}}{}_{\dot{\gamma}}.
\end{equation}

\paragraph{Consistency of Raising and Lowering (Eq. 4.57):}
Given $\psi^\alpha = \epsilon^{\alpha\beta}\psi_\beta$ and $\psi_\alpha = \epsilon_{\alpha\beta}\psi^\beta$ (and $\bar{\chi}^{\dot{\alpha}} = \epsilon^{\dot{\alpha}\dot{\beta}}\bar{\chi}_{\dot{\beta}}$, $\bar{\chi}_{\dot{\alpha}} = \epsilon_{\dot{\alpha}\dot{\beta}}\bar{\chi}^{\dot{\beta}}$), raising and then lowering returns the original spinor:
\begin{equation}
    \epsilon_{\alpha\beta}\psi^\beta = \epsilon_{\alpha\beta}\epsilon^{\beta\gamma}\psi_\gamma = \delta_\alpha{}^\gamma \psi_\gamma = \psi_\alpha, \qquad \epsilon_{\dot{\alpha}\dot{\beta}}\bar{\chi}^{\dot{\beta}} = \epsilon_{\dot{\alpha}\dot{\beta}}\epsilon^{\dot{\beta}\dot{\gamma}}\bar{\chi}_{\dot{\gamma}} = \delta_{\dot{\alpha}}{}^{\dot{\gamma}}\bar{\chi}_{\dot{\gamma}} = \bar{\chi}_{\dot{\alpha}}.
\end{equation}

\paragraph{Derivation of $SL(2,\mathbb{C})$ Invariance of $\epsilon_{\alpha\beta}$ (Eq. 4.59) and Dual Transformation (Eqs. 4.60 -- 4.68):}
Let an $SL(2,\mathbb{C})$ matrix have components $R_\alpha{}^\beta$ (where $\alpha$ is the row index and $\beta$ is the column index, as in Eqs.~4.48 and 4.60). In explicit index notation, $(R^T \epsilon R)_{\alpha\beta}$ is:
\begin{equation}
    R_\gamma{}^\alpha \epsilon_{\gamma\delta} R_\delta{}^\beta = R_1{}^\alpha \epsilon_{12} R_2{}^\beta + R_2{}^\alpha \epsilon_{21} R_1{}^\beta = R_1{}^\alpha R_2{}^\beta - R_2{}^\alpha R_1{}^\beta = (\det R)\epsilon_{\alpha\beta} = \epsilon_{\alpha\beta}.
\end{equation}
Similarly, the cofactor formula for the $2\times 2$ inverse of a unimodular matrix ($\det R = 1$) gives Eq.~(4.65):
\begin{equation}
    \epsilon^{\alpha\beta}R_\beta{}^\gamma \epsilon_{\gamma\delta} = (R^{-1})_\delta{}^\alpha.
\end{equation}
Substituting this into Eq.~(4.62):
\begin{equation}
    \psi'^\alpha = \epsilon^{\alpha\beta}\psi'_\beta = \epsilon^{\alpha\beta}R_\beta{}^\gamma \psi_\gamma = \epsilon^{\alpha\beta}R_\beta{}^\gamma \epsilon_{\gamma\delta}\psi^\delta = (R^{-1})_\delta{}^\alpha \psi^\delta = (R^{-1})_\beta{}^\alpha \psi^\beta,
\end{equation}
which proves Eq.~(4.66). Checking invariance of the spinor inner product $\psi^\alpha \chi_\alpha$ (Eqs.~4.67--4.68):
\begin{equation}
    \psi'^\alpha \chi'_\alpha = (R^{-1})_\beta{}^\alpha \psi^\beta R_\alpha{}^\gamma \chi_\gamma = (R^{-1})_\beta{}^\alpha R_\alpha{}^\gamma \psi^\beta \chi_\gamma = \delta_\beta{}^\gamma \psi^\beta \chi_\gamma = \psi^\beta \chi_\beta.
\end{equation}
For commuting spinors (Eq.~4.69):
\begin{equation}
    \psi^\alpha \chi_\alpha = \epsilon^{\alpha\beta}\psi_\beta \chi_\alpha = \epsilon^{\beta\alpha}\psi_\alpha \chi_\beta = -\epsilon^{\alpha\beta}\psi_\alpha \chi_\beta = -\psi_\alpha \chi^\alpha.
\end{equation}

\subsection{Four-Vectors as $2\times 2$ Hermitian Matrices and Spinor Bilinears (Eqs. 4.35 -- 4.46, 4.72 -- 4.95)}

\subsubsection{Part A: Derivation of Eqs. (4.35) -- (4.46)}
With $\eta_{\mu\nu} = \mathrm{diag}(-1,1,1,1)$, lowering the index on $x^\mu = (x^0, x^1, x^2, x^3)$ gives $x_\mu = \eta_{\mu\nu}x^\nu = (-x^0, x^1, x^2, x^3)$. With $(\sigma^\mu)_{\alpha\dot{\beta}} = (\mathbf{1}, \sigma^1, \sigma^2, \sigma^3)_{\alpha\dot{\beta}}$:
\begin{equation}
    X_{\alpha\dot{\beta}} = x_\mu (\sigma^\mu)_{\alpha\dot{\beta}} = -x^0 \delta_{\alpha\dot{\beta}} + x^i (\sigma_i)_{\alpha\dot{\beta}} = \begin{pmatrix} -x^0 + x^3 & x^1 - ix^2 \\ x^1 + ix^2 & -x^0 - x^3 \end{pmatrix}_{\alpha\dot{\beta}}.
\end{equation}
Computing the determinant (Eqs.~4.36 and 4.78):
\begin{equation}
    \det X = (-x^0 + x^3)(-x^0 - x^3) - (x^1 - ix^2)(x^1 + ix^2) = (x^0)^2 - (x^1)^2 - (x^2)^2 - (x^3)^2 = -x^\mu x_\mu.
\end{equation}
Under $X'_{\alpha\dot{\beta}} = R_\alpha{}^\gamma X_{\gamma\dot{\delta}} (R^\dagger)^{\dot{\delta}}{}_{\dot{\beta}}$ where $\det R = 1$, we have $\det X' = (\det R)(\det X)(\det R^\dagger) = \det X$, so $-(x')^\mu x'_\mu = -x^\mu x_\mu$. For a pure spatial rotation leaving $x^0$ invariant ($x'^0 = x^0$ for $x^i = 0$), setting $X_{\gamma\dot{\delta}} = -x^0 \delta_{\gamma\dot{\delta}}$ requires:
\begin{equation}
    R_\alpha{}^\gamma \delta_{\gamma\dot{\delta}} (R^\dagger)^{\dot{\delta}}{}_{\dot{\beta}} = \delta_{\alpha\dot{\beta}} \implies R R^\dagger = \mathbf{1} \implies R \in SU(2).
\end{equation}

\subsubsection{Part B: Derivation of Eqs. (4.75) -- (4.91) and Resolution of Typo in Eq. (4.88)}
Following the index convention of pages 10--11:
\begin{itemize}
    \item $\psi_L^\beta$ carries an upper undotted index and transforms as (Eq.~4.84):
    \begin{equation}
        \psi_L^\alpha \longrightarrow \psi'^\alpha_L = R^\alpha{}_\beta \psi_L^\beta.
    \end{equation}
    \item Its Hermitian conjugate $\psi_L^{\dagger\dot{\alpha}} \equiv (\psi_L^\alpha)^*$ carries an upper dotted index and transforms as (Eq.~4.85):
    \begin{equation}
        \psi_L^{\dagger\dot{\alpha}} \longrightarrow \psi'^{\dagger\dot{\alpha}}_L = (R^\alpha{}_\beta)^* (\psi_L^\beta)^* = \psi_L^{\dagger\dot{\beta}} (R^\dagger)_{\dot{\beta}}{}^{\dot{\alpha}},
    \end{equation}
    where $(R^\dagger)_{\dot{\beta}}{}^{\dot{\alpha}} \equiv (R^\alpha{}_\beta)^*$.
    \item In Eqs.~(4.80)--(4.81), $X_{\dot{\alpha}\beta} = x_\mu (\sigma^\mu)_{\dot{\alpha}\beta}$ transforms as:
    \begin{equation}
        X'_{\dot{\alpha}\beta} = x'_\mu (\sigma^\mu)_{\dot{\alpha}\beta} = (R^\dagger)_{\dot{\alpha}}{}^{\dot{\gamma}} X_{\dot{\gamma}\delta} R^\delta{}_\beta = x_\nu (R^\dagger)_{\dot{\alpha}}{}^{\dot{\gamma}} (\sigma^\nu)_{\dot{\gamma}\delta} R^\delta{}_\beta.
    \end{equation}
    \item Meanwhile, under $x'^\mu = \Lambda^\mu{}_\nu x^\nu$, the covector $x_\mu = \eta_{\mu\rho}x^\rho$ transforms with the inverse Lorentz matrix (Eq.~4.82):
    \begin{equation}
        x'_\mu = x_\nu (\Lambda^{-1})^\nu{}_\mu.
    \end{equation}
    Equating the two expressions for $X'_{\dot{\alpha}\beta}$ for arbitrary $x_\nu$ yields Eq.~(4.83):
    \begin{equation}
        (R^\dagger)_{\dot{\alpha}}{}^{\dot{\gamma}} (\sigma^\mu)_{\dot{\gamma}\delta} R^\delta{}_\beta = (\Lambda^{-1})^\mu{}_\nu (\sigma^\nu)_{\dot{\alpha}\beta}.
    \end{equation}
    \item Substituting this into the transformation of $V^\mu = \psi_L^{\dagger\dot{\alpha}} (\sigma^\mu)_{\dot{\alpha}\beta} \psi_L^\beta$ (Eqs.~4.86--4.88):
    \begin{align}
        V^\mu \longrightarrow V'^\mu &= \psi'^{\dagger\dot{\alpha}}_L (\sigma^\mu)_{\dot{\alpha}\beta} \psi'^\beta_L \nonumber \\
        &= \psi_L^{\dagger\dot{\gamma}} (R^\dagger)_{\dot{\gamma}}{}^{\dot{\alpha}} (\sigma^\mu)_{\dot{\alpha}\beta} R^\beta{}_\delta \psi_L^\delta \nonumber \\
        &= \psi_L^{\dagger\dot{\gamma}} \left[(\Lambda^{-1})^\mu{}_\nu (\sigma^\nu)_{\dot{\gamma}\delta}\right] \psi_L^\delta = (\Lambda^{-1})^\mu{}_\nu V^\nu.
    \end{align}
    \emph{Typo Note in Eq. (4.88):} Equation (4.88) in the lecture notes writes $\Lambda^\mu{}_\nu$ in the middle expression and $(\Lambda^{-1})^\mu{}_\nu$ on the right-hand side. If $R = \exp(-\frac{i}{2}\theta^i\sigma_i - \frac{1}{2}\eta^i\sigma_i)$ is the active transformation on $\psi_L$, then $V'^\mu = \Lambda^\mu{}_\nu V^\nu$ (as proven in Exercise 4.1 below), whereas if $\Lambda$ is defined via $X' = R^\dagger X R$, $\Lambda$ is replaced by $\Lambda^{-1}$.
\end{itemize}

\subsection{Solution to Exercise 4.1 and Exercise 4.2 (Eqs. 4.92 -- 4.95)}

\subsubsection{Exercise 4.1: Explicit Transformation of $V^\mu = (\psi_L^\dagger \psi_L, \psi_L^\dagger \sigma^i \psi_L)$ Under Boosts}
On the left-handed Weyl spinor $(\frac{1}{2},0)$, $B_i = \frac{1}{2}(j_i - ik_i) = 0$ and $A_i = \frac{1}{2}(j_i + ik_i) = \frac{1}{2}\sigma_i$:
\begin{equation}
    j_i = A_i + B_i = \frac{1}{2}\sigma_i, \qquad k_i = -i(A_i - B_i) = -\frac{i}{2}\sigma_i.
\end{equation}
Under a Lorentz boost with rapidity $\vec{\eta} = \eta \hat{n}$ ($\hat{n}^i \hat{n}^i = 1$), the $SL(2,\mathbb{C})$ matrix acting on $\psi_L^\alpha$ is:
\begin{equation}
    R^\alpha{}_\beta = \left(\exp\left(-i\eta^i k_i\right)\right)^\alpha{}_\beta = \left(\exp\left(-\frac{1}{2}\eta \hat{n}^i \sigma_i\right)\right)^\alpha{}_\beta = \cosh\left(\frac{\eta}{2}\right)\delta^\alpha{}_\beta - \sinh\left(\frac{\eta}{2}\right)\hat{n}^i (\sigma_i)^\alpha{}_\beta.
\end{equation}
Since the Pauli matrices are Hermitian, $((\sigma_i)^\alpha{}_\beta)^* = (\sigma_i)_{\dot{\beta}}{}^{\dot{\alpha}}$, the Hermitian conjugate matrix is:
\begin{equation}
    (R^\dagger)_{\dot{\gamma}}{}^{\dot{\alpha}} = \cosh\left(\frac{\eta}{2}\right)\delta_{\dot{\gamma}}{}^{\dot{\alpha}} - \sinh\left(\frac{\eta}{2}\right)\hat{n}^i (\sigma_i)_{\dot{\gamma}}{}^{\dot{\alpha}}.
\end{equation}
Using the Pauli matrix product identity in explicit index notation:
\begin{equation}
    (\sigma_i)_{\dot{\alpha}\beta}(\sigma_j)^\beta{}_\delta = \delta_{ij}\delta_{\dot{\alpha}\delta} + i\epsilon_{ijk}(\sigma_k)_{\dot{\alpha}\delta},
\end{equation}
we compute the transformations of $V^0$ and $V^k$:
\begin{enumerate}
    \item \textbf{Transformation of $V^0$:}
    \begin{align}
        (R^\dagger)_{\dot{\gamma}}{}^{\dot{\alpha}}\delta_{\dot{\alpha}\beta}R^\beta{}_\delta &= \left[\cosh\left(\frac{\eta}{2}\right)\delta_{\dot{\gamma}\beta} - \sinh\left(\frac{\eta}{2}\right)\hat{n}^i(\sigma_i)_{\dot{\gamma}\beta}\right]\left[\cosh\left(\frac{\eta}{2}\right)\delta^\beta{}_\delta - \sinh\left(\frac{\eta}{2}\right)\hat{n}^j(\sigma_j)^\beta{}_\delta\right] \nonumber \\
        &= \cosh^2\left(\frac{\eta}{2}\right)\delta_{\dot{\gamma}\delta} - 2\sinh\left(\frac{\eta}{2}\right)\cosh\left(\frac{\eta}{2}\right)\hat{n}^i(\sigma_i)_{\dot{\gamma}\delta} \nonumber \\
        &\quad + \sinh^2\left(\frac{\eta}{2}\right)\hat{n}^i\hat{n}^j\left(\delta_{ij}\delta_{\dot{\gamma}\delta} + i\epsilon_{ijk}(\sigma_k)_{\dot{\gamma}\delta}\right) \nonumber \\
        &= \cosh(\eta)\delta_{\dot{\gamma}\delta} - \sinh(\eta)\hat{n}^i(\sigma_i)_{\dot{\gamma}\delta}.
    \end{align}
    Contracting with $\psi_L^{\dagger\dot{\gamma}}$ and $\psi_L^\delta$ yields:
    \begin{equation}
        V'^0 = \cosh(\eta)V^0 - \sinh(\eta)\hat{n}^i V^i.
    \end{equation}

    \item \textbf{Transformation of $V^k$:}
    First compute $(\sigma^k)_{\dot{\alpha}\beta}R^\beta{}_\delta$:
    \begin{equation}
        (\sigma^k)_{\dot{\alpha}\beta}R^\beta{}_\delta = \cosh\left(\frac{\eta}{2}\right)(\sigma^k)_{\dot{\alpha}\delta} - \sinh\left(\frac{\eta}{2}\right)\hat{n}^j\left(\delta_{kj}\delta_{\dot{\alpha}\delta} + i\epsilon_{kjm}(\sigma_m)_{\dot{\alpha}\delta}\right).
    \end{equation}
    Now multiply on the left by $(R^\dagger)_{\dot{\gamma}}{}^{\dot{\alpha}}$:
    \begin{align}
        (R^\dagger)_{\dot{\gamma}}{}^{\dot{\alpha}}(\sigma^k)_{\dot{\alpha}\beta}R^\beta{}_\delta &= \cosh^2\left(\frac{\eta}{2}\right)(\sigma^k)_{\dot{\gamma}\delta} - \cosh\left(\frac{\eta}{2}\right)\sinh\left(\frac{\eta}{2}\right)\hat{n}^j\left(\delta_{kj}\delta_{\dot{\gamma}\delta} + i\epsilon_{kjm}(\sigma_m)_{\dot{\gamma}\delta}\right) \nonumber \\
        &\quad - \sinh\left(\frac{\eta}{2}\right)\cosh\left(\frac{\eta}{2}\right)\hat{n}^i\left(\delta_{ik}\delta_{\dot{\gamma}\delta} + i\epsilon_{ikm}(\sigma_m)_{\dot{\gamma}\delta}\right) \nonumber \\
        &\quad + \sinh^2\left(\frac{\eta}{2}\right)\hat{n}^i\hat{n}^j\left[\delta_{kj}(\sigma_i)_{\dot{\gamma}\delta} + i\epsilon_{kjm}\left(\delta_{im}\delta_{\dot{\gamma}\delta} + i\epsilon_{imn}(\sigma_n)_{\dot{\gamma}\delta}\right)\right] \nonumber \\
        &= (\sigma^k)_{\dot{\gamma}\delta} - \sinh(\eta)\hat{n}^k \delta_{\dot{\gamma}\delta} + (\cosh(\eta) - 1)\hat{n}^k \hat{n}^n (\sigma_n)_{\dot{\gamma}\delta}.
    \end{align}
    Contracting with $\psi_L^{\dagger\dot{\gamma}}$ and $\psi_L^\delta$ yields:
    \begin{equation}
        V'^k = V^k - \sinh(\eta)\hat{n}^k V^0 + (\cosh(\eta) - 1)\hat{n}^k (\hat{n}^n V^n).
    \end{equation}
\end{enumerate}

\subsubsection{Exercise 4.2: Explicit Transformation of $\bar{V}^\mu = \psi_R^\dagger \bar{\sigma}^\mu \psi_R$ and Origin of $\bar{\sigma}^0 = -\mathbf{1}$}
For a right-handed Weyl spinor $\psi_R$ transforming in $(0,\frac{1}{2})$, $A_i = \frac{1}{2}(j_i + ik_i) = 0$ and $B_i = \frac{1}{2}(j_i - ik_i) = \frac{1}{2}\sigma_i$, so:
\begin{equation}
    j_i = A_i + B_i = \frac{1}{2}\sigma_i, \qquad k_i = -i(A_i - B_i) = +\frac{i}{2}\sigma_i.
\end{equation}
Because $k_i$ has the opposite sign ($+\frac{i}{2}\sigma_i$) compared to $\psi_L$ ($-\frac{i}{2}\sigma_i$) while $j_i$ has the same sign, a boost by rapidity $\vec{\eta}$ acts on $\psi_R$ with $\eta \to -\eta$:
\begin{align}
    (\psi_R^\dagger \mathbf{1} \psi_R)' &= \cosh(\eta)(\psi_R^\dagger \mathbf{1} \psi_R) + \sinh(\eta)\hat{n}^i(\psi_R^\dagger \sigma^i \psi_R), \\
    (\psi_R^\dagger \sigma^k \psi_R)' &= (\psi_R^\dagger \sigma^k \psi_R) + \sinh(\eta)\hat{n}^k(\psi_R^\dagger \mathbf{1} \psi_R) + (\cosh(\eta)-1)\hat{n}^k \hat{n}^n(\psi_R^\dagger \sigma^n \psi_R).
\end{align}
Thus, for $\bar{V}^\mu = (\bar{V}^0, \bar{V}^k)$ to transform with the \textbf{same} boost matrix $\Lambda^\mu{}_\nu$ as $V^\mu$, $\bar{V}^0 = -\psi_R^\dagger \mathbf{1} \psi_R$ must have a relative minus sign compared to $\bar{V}^k = +\psi_R^\dagger \sigma^k \psi_R$, i.e., $\bar{\sigma}^\mu = (-\mathbf{1}, \sigma^i)$.

\paragraph{Index/Metric Origin of the Minus Sign in $\bar{\sigma}^0 = -\mathbf{1}$:}
In spinor index notation with $\epsilon^{\dot{\alpha}\dot{\gamma}} = -\epsilon^{\alpha\gamma}$, raising both indices of $(\sigma^\mu)_{\dot{\gamma}\delta}$ via:
\begin{equation}
    (\bar{\sigma}^\mu)^{\beta\dot{\alpha}} = -\epsilon^{\beta\delta}\epsilon^{\dot{\alpha}\dot{\gamma}}(\sigma^\mu)_{\dot{\gamma}\delta}
\end{equation}
gives:
\begin{itemize}
    \item For $\mu=0$, $(\sigma^0)_{\dot{\gamma}\delta} = \delta_{\dot{\gamma}\delta}$, so $-\epsilon^{\beta\delta}\epsilon^{\dot{\alpha}\dot{\gamma}}\delta_{\dot{\gamma}\delta} = -\epsilon^{\beta\delta}(-\epsilon^{\dot{\alpha}\delta}) = -\delta^{\beta\dot{\alpha}} = -(\mathbf{1})^{\beta\dot{\alpha}}$.
    \item For $\mu=i$, $-\epsilon^{\beta\delta}\epsilon^{\dot{\alpha}\dot{\gamma}}(\sigma^i)_{\dot{\gamma}\delta} = +(\sigma^{i*})^{\beta\dot{\alpha}}$.
\end{itemize}
Equivalently, from the Clifford identity $(\sigma^\mu)_{\dot{\alpha}\beta}(\bar{\sigma}^\nu)^{\beta\dot{\gamma}} + (\sigma^\nu)_{\dot{\alpha}\beta}(\bar{\sigma}^\mu)^{\beta\dot{\gamma}} = 2\eta^{\mu\nu}\delta_{\dot{\alpha}}{}^{\dot{\gamma}}$ in $(-+++)$ signature, $\eta^{00} = -1$ forces $\sigma^0 \bar{\sigma}^0 = -\mathbf{1} \implies \bar{\sigma}^0 = -\mathbf{1}$, while $\eta^{ii} = +1$ forces $\sigma^i \bar{\sigma}^i = +\mathbf{1} \implies \bar{\sigma}^i = +\sigma^i$.

\vspace{0.5em}\hrule\vspace{1em}

\section{One-Particle States, Little Group, and Casimir Operators (Section 4.1.3)}

\subsection{First Casimir Operator $C_1 = -P_\mu P^\mu$ (Eqs. 4.96 -- 4.121)}
Using $[P^\mu, P^\nu] = 0$ and $[L^{\mu\nu}, P^\rho] = i(\eta^{\nu\rho}P^\mu - \eta^{\mu\rho}P^\nu)$:
\begin{align}
    [L^{\mu\nu}, -P_\rho P^\rho] &= -\eta_{\rho\sigma}\left([L^{\mu\nu}, P^\rho]P^\sigma + P^\rho[L^{\mu\nu}, P^\sigma]\right) \nonumber \\
    &= -i\eta_{\rho\sigma}\left((\eta^{\nu\rho}P^\mu - \eta^{\mu\rho}P^\nu)P^\sigma + P^\rho(\eta^{\nu\sigma}P^\mu - \eta^{\mu\sigma}P^\nu)\right) \nonumber \\
    &= -i\left(P^\mu P^\nu - P^\nu P^\mu + P^\nu P^\mu - P^\mu P^\nu\right) = 0.
\end{align}

\subsection{Induced Representations and Wigner Rotation (Eqs. 4.122 -- 4.128)}
For reference momentum $k^\mu$ and boost $L(p)^\mu{}_\nu k^\nu = p^\mu$, let $W(\Lambda, p)^\mu{}_\nu \equiv (L^{-1}(\Lambda p))^\mu{}_\rho \Lambda^\rho{}_\sigma L(p)^\sigma{}_\nu$. Acting on $k^\nu$ in explicit index notation (Eq.~4.126):
\begin{align}
    W(\Lambda, p)^\mu{}_\nu k^\nu &= (L^{-1}(\Lambda p))^\mu{}_\rho \Lambda^\rho{}_\sigma \left(L(p)^\sigma{}_\nu k^\nu\right) \nonumber \\
    &= (L^{-1}(\Lambda p))^\mu{}_\rho \left(\Lambda^\rho{}_\sigma p^\sigma\right) = (L^{-1}(\Lambda p))^\mu{}_\rho (\Lambda p)^\rho = k^\mu,
\end{align}
so $W(\Lambda, p) \in G_k$ (Eq.~4.127), and:
\begin{equation}
    U(\Lambda)|p,\lambda\rangle = U(L(\Lambda p))U(W(\Lambda,p))|k,\lambda\rangle = \sum_{\lambda'} |\Lambda p, \lambda'\rangle D_{\lambda'\lambda}(W(\Lambda,p)).
\end{equation}

\subsection{The Pauli-Lubanski Vector and Second Casimir $C_2 = W_\mu W^\mu$ (Eqs. 4.129 -- 4.136)}
The Pauli-Lubanski vector is defined as (Eq.~4.129):
\begin{equation}
    W^\mu = \frac{1}{2}\epsilon^{\mu\nu\rho\sigma}P_\nu L_{\rho\sigma}.
\end{equation}
\begin{enumerate}
    \item \textbf{Proof that $W^\mu P_\mu = 0$ and $[W^\mu, P^\alpha] = 0$ (Eq. 4.130):}
    \begin{align}
        [W^\mu, P^\alpha] &= \frac{1}{2}\epsilon^{\mu\nu\rho\sigma}P_\nu [L_{\rho\sigma}, P^\alpha] = \frac{i}{2}\epsilon^{\mu\nu\rho\sigma}P_\nu \left(\delta_\sigma{}^\alpha P_\rho - \delta_\rho{}^\alpha P_\sigma\right) \nonumber \\
        &= \frac{i}{2}\left(\epsilon^{\mu\nu\rho\alpha}P_\nu P_\rho - \epsilon^{\mu\nu\alpha\sigma}P_\nu P_\sigma\right) = 0, \\[4pt]
        W^\mu P_\mu &= \frac{1}{2}\epsilon^{\mu\nu\rho\sigma}P_\mu P_\nu L_{\rho\sigma} = 0.
    \end{align}
    \item \textbf{Proof that $[W^\lambda, L^{\mu\nu}] = i(W^\mu \eta^{\lambda\nu} - W^\nu \eta^{\lambda\mu})$ (Eq. 4.131):}
    Using the 5-index antisymmetrization identity $\epsilon^{[\lambda\tau\beta\gamma}\eta^{\nu]\mu} = 0$ and the Leibniz rule:
    \begin{align}
        [L^{\mu\nu}, W^\lambda] &= \frac{1}{2}\epsilon^{\lambda\alpha\beta\gamma}\left([L^{\mu\nu}, P_\alpha]L_{\beta\gamma} + P_\alpha [L^{\mu\nu}, L_{\beta\gamma}]\right) \nonumber \\
        &= \frac{i}{2}\left[\epsilon^{\lambda\nu\beta\gamma}P^\mu L_{\beta\gamma} + \epsilon^{\lambda\alpha\nu\gamma}P_\alpha L^\mu{}_\gamma + \epsilon^{\lambda\alpha\beta\nu}P_\alpha L_\beta{}^\mu - (\mu \leftrightarrow \nu)\right] \nonumber \\
        &= i\left(\eta^{\lambda\nu}W^\mu - \eta^{\lambda\mu}W^\nu\right) \iff [W^\lambda, L^{\mu\nu}] = i\left(W^\mu \eta^{\lambda\nu} - W^\nu \eta^{\lambda\mu}\right).
    \end{align}
    \item \textbf{Proof that $C_2 = W_\mu W^\mu$ is a Casimir Operator (Eq. 4.132):}
    \begin{align}
        [L^{\mu\nu}, W_\lambda W^\lambda] &= -i\eta_{\lambda\tau}\left((W^\mu\eta^{\lambda\nu} - W^\nu\eta^{\lambda\mu})W^\tau + W^\lambda(W^\mu\eta^{\tau\nu} - W^\nu\eta^{\tau\mu})\right) \nonumber \\
        &= -i\left(W^\mu W^\nu - W^\nu W^\mu + W^\nu W^\mu - W^\mu W^\nu\right) = 0.
    \end{align}
    \item \textbf{Rest-Frame Evaluation of $W^\mu$ and $C_2$ (Eqs. 4.133 -- 4.136):}
    In the rest frame $P^\mu = (m,0,0,0) \implies P_0 = -m, P_i = 0$. With $\epsilon^{0123} = -1$ ($\epsilon^{0ijk} = -\epsilon_{ijk}$) and $J_i = -\frac{1}{2}\epsilon_{ijk}L^{jk}$:
    \begin{equation}
        W^0 = 0, \qquad W^i = \frac{1}{2}\epsilon^{i0jk}P_0 L_{jk} = \frac{1}{2}\left(-\epsilon^{0ijk}\right)(-m)\left(-\epsilon_{jkl}J_l\right) = m J_i,
    \end{equation}
    so $C_2 = \eta_{\mu\nu}W^\mu W^\nu = -(W^0)^2 + \delta_{ij}W^i W^j = m^2 \vec{J}^2 = m^2 s(s+1)$.
\end{enumerate}

\subsection{Massless Little Group $ISO(2)$ and Helicity (Eqs. 4.108 -- 4.115, 4.137 -- 4.149)}
For the standard null reference momentum $k^\mu = (E, 0, 0, E)$, $k_\mu = (-E, 0, 0, E)$.
\begin{enumerate}
    \item \textbf{Little Group Generators:}
    Requiring $\omega_{\rho\sigma}(l^{\rho\sigma})^\mu{}_\nu k^\nu = 2\omega^\mu{}_\sigma k^\sigma = 0$ with $\omega_{0i} = -\eta_i$ and $\omega_{ij} = \epsilon_{ijk}\theta_k$ gives $\eta_3 = 0$, $\theta_2 = -\eta_1$, and $\theta_1 = \eta_2$, so $\theta_i J_i + \eta_i K_i = \theta_3 J_3 + \eta_1(K_1 - J_2) + \eta_2(K_2 + J_1)$, spanned by:
    \begin{equation}
        J_3, \qquad N_1 = K_1 - J_2 = -\Pi_2, \qquad N_2 = K_2 + J_1 = \Pi_1.
    \end{equation}
    \item \textbf{Commutation Relations (Eqs. 4.111, 4.139, 4.140):}
    \begin{align}
        [J_3, \Pi_1] &= [J_3, J_1 + K_2] = iJ_2 - iK_1 = i\Pi_2, \\
        [J_3, \Pi_2] &= [J_3, J_2 - K_1] = -iJ_1 - iK_2 = -i\Pi_1, \\
        [\Pi_1, \Pi_2] &= [J_1 + K_2, J_2 - K_1] = iJ_3 - iJ_3 = 0.
    \end{align}
    \item \textbf{Pauli-Lubanski Vector on Massless States (Eqs. 4.138, 4.145 -- 4.147):}
    Evaluating $W^\mu = \frac{1}{2}\epsilon^{\mu\nu\rho\sigma}P_\nu L_{\rho\sigma}$ on $P_0 = -E, P_3 = E, P_1 = P_2 = 0$ (with $L_{ij} = -\epsilon_{ijk}J_k$ and $L_{0i} = K_i$):
    \begin{equation}
        W^0 = E J_3, \qquad W^3 = E J_3, \qquad W^1 = -E(J_1 + K_2) = -E\Pi_1, \qquad W^2 = -E(J_2 - K_1) = -E\Pi_2.
    \end{equation}
    When $\Pi_1 = \Pi_2 = 0$ and $J_3|k,h\rangle = h|k,h\rangle$, we obtain $W^\mu = hP^\mu$ and $W_\mu W^\mu = h^2 P_\mu P^\mu = 0$.
\end{enumerate}

\vspace{0.5em}\hrule\vspace{1em}

\section{Complete Solutions to Problems 4.1 and 4.2 (Section 4.2)}

\subsection{Problem 4.1: Decomposition of Rank-Two Lorentz Tensors and Massive Spin-2 Polarizations}

\subsubsection*{(a) Decomposition of $(\frac{1}{2}, \frac{1}{2}) \otimes (\frac{1}{2}, \frac{1}{2})$}
Using the $SU(2)$ Clebsch-Gordan decomposition $\frac{1}{2} \otimes \frac{1}{2} = 0 \oplus 1$ on each chiral factor:
\begin{equation}
    \left(\frac{1}{2}, \frac{1}{2}\right) \otimes \left(\frac{1}{2}, \frac{1}{2}\right) = \left(\frac{1}{2} \otimes \frac{1}{2}, \frac{1}{2} \otimes \frac{1}{2}\right) = (0 \oplus 1, 0 \oplus 1) = (0,0) \oplus (1,0) \oplus (0,1) \oplus (1,1).
\end{equation}
Dimension check using $\dim(j_L, j_R) = (2j_L+1)(2j_R+1)$:
\begin{equation}
    4 \times 4 = 16 = 1 + 3 + 3 + 9.
\end{equation}

\subsubsection*{(b) Explicit Index Decomposition of $T^{\mu\nu}$}
Define $T^{[\mu\nu]} = \frac{1}{2}(T^{\mu\nu} - T^{\nu\mu})$ and $T^{(\mu\nu)} = \frac{1}{2}(T^{\mu\nu} + T^{\nu\mu})$. Then:
\begin{equation}
    T^{\mu\nu} = \underbrace{\frac{1}{4}\eta^{\mu\nu}T^\rho{}_\rho}_{(0,0),\ \dim=1} + \underbrace{T^{[\mu\nu]}}_{(1,0)\oplus(0,1),\ \dim=3+3=6} + \underbrace{\left(T^{(\mu\nu)} - \frac{1}{4}\eta^{\mu\nu}T^\rho{}_\rho\right)}_{(1,1),\ \dim=9}.
\end{equation}

\subsubsection*{(c) Restriction to the Spatial Rotation Subgroup $SU(2)_{\mathrm{rot}}$}
Because $j_i = A_i + B_i$ generates the diagonal subgroup $SU(2)_{\mathrm{rot}} \subset SU(2)_L \times SU(2)_R$, the representation $(j_L, j_R)$ decomposes under spatial rotations via angular momentum addition:
\begin{equation}
    (j_L, j_R)\big|_{SU(2)_{\mathrm{rot}}} \simeq \bigoplus_{j=|j_L - j_R|}^{j_L + j_R} j.
\end{equation}
\begin{itemize}
    \item $(\frac{1}{2}, \frac{1}{2})\big|_{SU(2)_{\mathrm{rot}}} = 0 \oplus 1$ (1 time component $V^0$ with $j=0$, and 3 spatial components $V^i$ with $j=1$).
    \item $(1,0)\big|_{SU(2)_{\mathrm{rot}}} = 1$ and $(0,1)\big|_{SU(2)_{\mathrm{rot}}} = 1$ (self-dual and anti-self-dual combinations $\vec{E} \pm i\vec{B}$ each transform as 3-vectors of spin $j=1$).
    \item $(1,1)\big|_{SU(2)_{\mathrm{rot}}} = 0 \oplus 1 \oplus 2$ ($S^{00}$ has $j=0$ [1 component], $S^{0i}$ has $j=1$ [3 components], and traceless symmetric $S^{ij} - \frac{1}{3}\delta^{ij}S^{kk}$ has $j=2$ [5 components]).
\end{itemize}

\subsubsection*{(d) Rest-Frame Analysis of Massive Spin-2 Polarization Tensor $\epsilon^{\mu\nu}(p)$}
Given $\epsilon^{\mu\nu} = \epsilon^{\nu\mu}$, $\eta_{\mu\nu}\epsilon^{\mu\nu} = 0$, and $p_\mu \epsilon^{\mu\nu} = 0$ in the rest frame $p^\mu = (m,0,0,0) \implies p_\mu = (-m,0,0,0)$:
\begin{equation}
    p_\mu \epsilon^{\mu\nu} = -m\epsilon^{0\nu} = 0 \implies \epsilon^{0\nu} = \epsilon^{\nu 0} = 0 \quad (\forall \nu \in \{0,1,2,3\}).
\end{equation}
Only the $3\times 3$ spatial submatrix $\epsilon^{ij}$ ($i,j \in \{1,2,3\}$) survives, subject to:
\begin{equation}
    \epsilon^{ij} = \epsilon^{ji} \quad (6\text{ components}), \qquad \eta_{\mu\nu}\epsilon^{\mu\nu} = -\epsilon^{00} + \delta_{ij}\epsilon^{ij} = \delta_{ij}\epsilon^{ij} = 0 \quad (-1\text{ constraint}),
\end{equation}
leaving $6 - 1 = 5$ independent components.

\subsubsection*{(e) Spin-2 Irreducible Representation Under $SO(3)$}
Under the rest-frame little group $SO(3)$, a general spatial tensor $t^{ij}$ in $\mathbf{3} \otimes \mathbf{3} = \mathbf{1} \oplus \mathbf{3} \oplus \mathbf{5}$ ($1 \otimes 1 = 0 \oplus 1 \oplus 2$) decomposes into trace ($\mathbf{1}$, $j=0$), antisymmetric ($\mathbf{3}$, $j=1$), and symmetric traceless ($\mathbf{5}$, $j=2$) parts. Since $\epsilon^{ij}$ is symmetric and traceless, it transforms irreducibly in the $2(2)+1 = 5$-dimensional spin-$2$ ($j=2$) representation of $SO(3)$, describing the 5 polarization states of a massive spin-2 particle.

\vspace{0.5em}\hrule\vspace{1em}

\subsection{Problem 4.2: Explicit Index Proof of the Fierz Identity (Eqs. 4.150 -- 4.151)}

\paragraph{Step 1: Completeness Relation for $(\sigma^\mu)_{\alpha\dot{\alpha}} (\sigma_\mu)_{\beta\dot{\beta}}$ in $(-+++)$ Signature with $\epsilon_{\dot{\alpha}\dot{\beta}} = -\epsilon_{\alpha\beta}$}
Using $\eta_{\mu\nu} = \mathrm{diag}(-1, +1, +1, +1)$ and $\sigma^\mu = (\mathbf{1}, \sigma^1, \sigma^2, \sigma^3)$:
\begin{equation}
    (\sigma^\mu)_{\alpha\dot{\alpha}}(\sigma_\mu)_{\beta\dot{\beta}} = -(\sigma^0)_{\alpha\dot{\alpha}}(\sigma^0)_{\beta\dot{\beta}} + (\sigma^1)_{\alpha\dot{\alpha}}(\sigma^1)_{\beta\dot{\beta}} + (\sigma^2)_{\alpha\dot{\alpha}}(\sigma^2)_{\beta\dot{\beta}} + (\sigma^3)_{\alpha\dot{\alpha}}(\sigma^3)_{\beta\dot{\beta}}.
\end{equation}
Checking components against $\epsilon_{12} = +1, \epsilon_{21} = -1$ and the professor's convention $\epsilon_{\dot{1}\dot{2}} = -1, \epsilon_{\dot{2}\dot{1}} = +1$:
\begin{itemize}
    \item $(\alpha,\dot{\alpha},\beta,\dot{\beta}) = (1,1,1,1)$: $-1 + 0 + 0 + 1 = 0 = +2\epsilon_{11}\epsilon_{\dot{1}\dot{1}}$.
    \item $(\alpha,\dot{\alpha},\beta,\dot{\beta}) = (1,1,2,2)$: $-1 + 0 + 0 - 1 = -2 = +2\epsilon_{12}\epsilon_{\dot{1}\dot{2}}$ (since $\epsilon_{12}\epsilon_{\dot{1}\dot{2}} = (+1)(-1) = -1$).
    \item $(\alpha,\dot{\alpha},\beta,\dot{\beta}) = (1,2,2,1)$: $-0 + 1 + (-i)(i) + 0 = +2 = +2\epsilon_{12}\epsilon_{\dot{2}\dot{1}}$ (since $\epsilon_{12}\epsilon_{\dot{2}\dot{1}} = (+1)(+1) = +1$).
    \item $(\alpha,\dot{\alpha},\beta,\dot{\beta}) = (1,2,1,2)$: $-0 + 1 + (-i)^2 + 0 = 0 = +2\epsilon_{11}\epsilon_{\dot{2}\dot{2}}$.
\end{itemize}
By symmetry under $1 \leftrightarrow 2$, this proves the completeness identity for all index values:
\begin{equation}
    (\sigma^\mu)_{\alpha\dot{\alpha}}(\sigma_\mu)_{\beta\dot{\beta}} = +2\epsilon_{\alpha\beta}\epsilon_{\dot{\alpha}\dot{\beta}}.
\end{equation}

\paragraph{Step 2: Contraction with Spinors}
With $\chi\sigma^\mu\bar{\psi} = \chi^\alpha (\sigma^\mu)_{\alpha\dot{\alpha}}\bar{\psi}^{\dot{\alpha}}$ and $\lambda\sigma_\mu\bar{\rho} = \lambda^\beta (\sigma_\mu)_{\beta\dot{\beta}}\bar{\rho}^{\dot{\beta}}$, and using $\lambda_\alpha = \epsilon_{\alpha\beta}\lambda^\beta$ and $\bar{\psi}_{\dot{\beta}} = \epsilon_{\dot{\beta}\dot{\alpha}}\bar{\psi}^{\dot{\alpha}} = -\epsilon_{\dot{\alpha}\dot{\beta}}\bar{\psi}^{\dot{\alpha}}$:
\begin{enumerate}
    \item \textbf{For Commuting ($c$-number) Spinors ($\bar{\psi}^{\dot{\alpha}}\lambda^\beta = +\lambda^\beta \bar{\psi}^{\dot{\alpha}}$):}
    \begin{align}
        (\chi \sigma^\mu \bar{\psi})(\lambda \sigma_\mu \bar{\rho}) &= (\sigma^\mu)_{\alpha\dot{\alpha}}(\sigma_\mu)_{\beta\dot{\beta}} \chi^\alpha \bar{\psi}^{\dot{\alpha}} \lambda^\beta \bar{\rho}^{\dot{\beta}} \nonumber \\
        &= +2\epsilon_{\alpha\beta}\epsilon_{\dot{\alpha}\dot{\beta}} \chi^\alpha \lambda^\beta \bar{\psi}^{\dot{\alpha}} \bar{\rho}^{\dot{\beta}} \nonumber \\
        &= +2\left(\chi^\alpha \epsilon_{\alpha\beta}\lambda^\beta\right)\left(-\epsilon_{\dot{\beta}\dot{\alpha}}\bar{\psi}^{\dot{\alpha}}\bar{\rho}^{\dot{\beta}}\right) \nonumber \\
        &= -2\left(\chi^\alpha \lambda_\alpha\right)\left(\bar{\psi}_{\dot{\beta}}\bar{\rho}^{\dot{\beta}}\right) = -2(\chi\lambda)(\bar{\psi}\bar{\rho}).
    \end{align}
    \item \textbf{For Anticommuting (Grassmann) Spinors ($\bar{\psi}^{\dot{\alpha}}\lambda^\beta = -\lambda^\beta \bar{\psi}^{\dot{\alpha}}$):}
    Reordering $\bar{\psi}^{\dot{\alpha}}$ past $\lambda^\beta$ introduces an extra minus sign, yielding $+2(\chi\lambda)(\bar{\psi}\bar{\rho})$ (or $-2(\chi\lambda)(\bar{\psi}\bar{\rho})$ if the dotted contraction is defined as $\bar{\psi}\bar{\rho} = \bar{\psi}^{\dot{\alpha}}\bar{\rho}_{\dot{\alpha}}$).
\end{enumerate}

\end{document}

